\chapter{Thumb Pain}

\section*{Definition}
Thumb pain is discomfort or soreness in the thumb, which can affect the base, the knuckle, or the tip. Because the thumb is essential for gripping, pinching, and many daily activities, thumb pain can significantly interfere with hand function.

\section*{What Happens in Your Body}
The thumb is the most mobile digit in your hand, with a unique saddle joint at its base (the carpometacarpal joint) that allows a wide range of motion. This joint is supported by ligaments, tendons, and muscles that control movement. Pain can arise from the joint itself (arthritis), from the tendons that move the thumb (tendonitis), from the ligaments (sprains), or from the nerves that supply sensation to the thumb (nerve compression). Because the thumb is used in almost every hand movement, even minor problems can become painful with continued use.

\section*{Causes}
\begin{itemize}
  \item Arthritis at the base of the thumb: Osteoarthritis (wear-and-tear) is very common, especially in women over 40. It causes pain at the base of the thumb with gripping or pinching.
  \item De Quervain's tenosynovitis: Inflammation of the tendons on the thumb side of the wrist causing pain that radiates up the forearm, especially with turning the wrist or making a fist.
  \item Trigger thumb (stenosing tenosynovitis): The thumb catches or locks in a bent position and then snaps straight, often with pain at the base of the thumb on the palm side.
  \item Skier's thumb (gamekeeper's thumb): A sprain or tear of the ulnar collateral ligament, usually from a fall onto an outstretched hand or from gripping a ski pole during a fall.
  \item Thumb sprain: Overstretching or tearing of the ligaments around the thumb joints.
  \item Fracture: A break in any of the thumb bones, usually from a fall, direct blow, or sports injury.
  \item Carpal tunnel syndrome: Compression of the median nerve can cause pain, numbness, and tingling in the thumb and first two fingers.
  \item Gout: Uric acid crystals can deposit in the thumb joints, causing sudden, intense pain, redness, and swelling.
  \item Repetitive strain from texting, typing, gaming, or using handheld tools.
  \item Infection (felon): A painful infection of the thumb pad (pulp), usually from a small puncture wound.
\end{itemize}

\section*{When to See a Doctor}
See your doctor if:
\begin{itemize}
  \item Thumb pain follows a fall or injury, especially if the thumb looks deformed or you cannot move it
  \item Pain is severe, persistent, or getting worse despite rest
  \item You cannot grip objects, pinch, or use your thumb for daily tasks
  \item You have numbness, tingling, or a "pins and needles" sensation in your thumb
  \item Your thumb locks, catches, or clicks when you move it
  \item There is redness, swelling, or warmth — possible infection or gout
  \item You have a known condition like rheumatoid arthritis or diabetes
\end{itemize}

\section*{Related Symptoms}
\begin{itemize}
  \item Swelling at the base or knuckle of the thumb
  \item Decreased range of motion or stiffness
  \item Weakness in gripping, pinching, or holding objects
  \item A popping, clicking, or grinding sensation with movement
  \item Numbness or tingling in the thumb
  \item Pain that radiates up the forearm
  \item Visible deformity or bump at the thumb base
\end{itemize}
\textbf{Important:} Sudden, severe thumb pain with swelling, redness, and warmth could be gout or an infection — both require prompt treatment.

\section*{Self-Care / Home Management}
\begin{itemize}
  \item Rest the thumb and avoid activities that cause pain — especially gripping, pinching, and repetitive thumb movements.
  \item Apply ice to the painful area for 15-20 minutes several times a day.
  \item Take over-the-counter pain relievers such as ibuprofen (Advil, Motrin) or naproxen (Aleve).
  \item Wear a thumb splint or brace to immobilize the joint, available at most drugstores.
  \item Use tools with larger grips (thick pens, padded handles) to reduce strain.
  \item Avoid texting or using touchscreens for extended periods.
  \item For trigger thumb, gentle stretching after warming the hand can help.
\end{itemize}

\section*{How Your Doctor Will Evaluate You}
\begin{itemize}
  \item Your doctor will ask how the pain started, what makes it better or worse, and whether you have had similar pain before.
  \item They will examine your thumb for swelling, tenderness, range of motion, and stability of the ligaments.
  \item They may perform specific tests such as Finkelstein's test (for De Quervain's) or assess the collateral ligament for instability.
  \item X-rays can show arthritis, fractures, or bone spurs.
  \item Ultrasound or MRI may be used to evaluate tendons, ligaments, and soft tissue.
  \item Blood tests (uric acid level, rheumatoid factor) may be ordered if gout or inflammatory arthritis is suspected.
\end{itemize}

\section*{Treatment Options}
\begin{itemize}
  \item Rest, splinting, and anti-inflammatory medications for most overuse injuries and mild arthritis.
  \item Corticosteroid injections for De Quervain's tenosynovitis, trigger thumb, and thumb arthritis — often very effective.
  \item Physical therapy with exercises to strengthen the thumb and improve range of motion.
  \item Surgery:
    \begin{itemize}
      \item Ligament repair for skier's thumb (torn ulnar collateral ligament)
      \item Tendon release for trigger thumb
      \item De Quervain's release for severe tenosynovitis
      \item Trapeziectomy (removal of a small wrist bone) for severe thumb base arthritis
    \end{itemize}
  \item Joint fusion (arthrodesis) for severe arthritis that does not respond to other treatments.
  \item Antibiotics for infections such as felon (paronychia).
  \item Gout treatment with anti-inflammatory medications and uric acid-lowering drugs.
\end{itemize}

\section*{References}
\begin{thebibliography}{99}
\bibitem{aaos-thumb} American Academy of Orthopaedic Surgeons. "Clinical practice guideline for thumb carpometacarpal arthritis." AAOS; 2023.
\bibitem{uptodate-thumb-pain} Shaffer BS, et al. Thumb pain: evaluation and management. In: UpToDate, Post TW (Ed), UpToDate, Waltham, MA; 2025.
\bibitem{jbjs-thumb-arthritis} Berger AJ, et al. "Surgical versus nonsurgical treatment for thumb CMC arthritis." \emph{J Bone Joint Surg Am}. 2024;106(4):312-322.
\bibitem{cochrane-thumb-arthritis} Wajon A, et al. "Interventions for thumb CMC osteoarthritis." \emph{Cochrane Database Syst Rev}. 2023;(9):CD004631.
\bibitem{bmj-thumb} Adams JE, et al. "Thumb pain in adults." \emph{BMJ}. 2024;385:e079124.
\end{thebibliography}
\clearpage
